%% LyX 2.5.1 created this file.  For more info, see https://www.lyx.org/.
%% Do not edit unless you really know what you are doing.
\documentclass[brazil]{article}
\usepackage[T1]{fontenc}
\usepackage[utf8]{luainputenc}
\usepackage{amsmath}
\usepackage{babel}
\begin{document}

\appendix

\section{Questões matemáticas e computacionais}

\subsection{Lei de Pareto como lei estatística}

Além da construção do modelo generativo que discutimos com mais detalhes previamente, outra atividade necessária na investigação de sistemas complexos adotando a abordagem proposta consiste na construção de ao menos uma lei estatística macroscópica que deve emergir deste modelo generativo a ser construído. Nesta pesquisa, não nos ocupamos desta parte da investigação; porém, caso o leitor esteja interessado, o livro elaborado por Altmann dedica um capítulo a esta questão, isto é, Como leis estatísticas podem (e devem) ser obtidas a partir dos dados disponíveis\cite{Altmann_2024}. 

Segundo Altmann\cite{Altmann_2024}, a lei de desigualdade de Pareto é a mais influente e provavelmente o primeiro exemplo de lei estatística como uma distribuição de lei de potência. Ela surgiu de uma análise empírica da distribuição de renda de diferentes países, buscando identificar qual a proporção $N\left(x\right)$ da população possui uma renda maior que $x$. Os dados que foram compilados por Pareto de fontes originais são principalmente do século XIX, e o resultado de sua análise foi proposto da seguinte forma: 
\[
\ln N\left(x\right)=\ln A-\tilde{\gamma}\ln\left(x\right)
\]
Com o mesmo valor $\tilde{\gamma}\approx1.5$ sendo observado em cenários completamente diferentes, Pareto também percebeu que essa aproximação lembra uma função de distribuição cumulativa complementar (\textit{Complementary cumulative distribution function} - CCDF). A alegação de que esta distribuição descreve a renda (e riqueza) em diferentes países e regiões é chamada de lei de Pareto.

O quão bem os dados se ajustam a esta lei é alvo de debate. Além de uma lei de potências, também foram feitas propostas de que os dados seriam bem descritos por distribuições do tipo Lognormal (ou Lognormal generalizado) ou pela lei exponencial de Boltzmann-Gibbs. Pesquisas empíricas mais recentes tem indicado que a distribuição de renda pode ser separada em uma região de baixa renda e uma região de maior renda, passando de uma distribuição exponencial para uma distribuição de lei de potência. De maneira geral, pesquisas recentes examinando os dados fornecidos pela World Distribution of Income tem sugerido uma distribuição bimodal \cite{diss_mods_00007721}.

De toda forma, mais de um século depois de ter sido proposta, a utilidade das distribuições do tipo Pareto para caracterizar distribuições de renda e riqueza continua a ser reconhecida nas análises econômicas modernas. Uma vez que estamos em posse da lei estatística, um caminho possível é procurar a explicação da origem da lei, uma tarefa que o próprio Pareto se dedicou inicialmente. Esta tarefa continua a ser desenvolvida até os dias de hoje, uma vez que nenhuma explicação foi capaz de fazer com que os cientistas considerassem este um caso encerrado.

O outro caminho possível é analisar suas consequências. Sobre isso, um debate que tem sido realizado é como a variação ou a falta de variação do exponente $\tilde{\gamma}$ afetaria o bem estar e a desigualdade na população; assim como foram realizadas discussões econômicas sobre como melhorá-lo, assumindo a validade da lei. Evidentemente, todas essas discussões estão imersas em controvérsias.

\subsubsection{Definição matemática}

Mas por hora, nos interessa que a lei de Pareto nos dá a probabilidade de uma pessoa ter exatamente o valor $x$ da variável de interesse, no nosso caso, a probabilidade de um agente ter uma riqueza no valor $w$ . A função de densidade de probabilidade (\textit{Probability Density Function} - PDF)\footnote{Esta função é usada para especificar a probabilidade de uma variável aleatória estar dentro de um determinado intervalo de valores.} conhecida como lei de Pareto, é \cite{pedroso}:

\[
p\left(w\right)=\frac{\alpha\beta^{\alpha}}{w^{\alpha+1}}
\]

Onde $\alpha>0$ é o coeficiente de Pareto (parâmetro de forma) que modela a distribuição e $\beta>0$ é a riqueza mínima (parâmetro de escala). O parâmetro que normalmente domina as discussões sobre a lei de Pareto refere-se ao índice de Pareto. Dessa forma, um objetivo comum é calcular o coeficiente de Pareto nas distribuições que estamos analisando. Para isso, iremos trabalhar com a função de distribuição cumulativa complementar.

A função de distribuição cumulativa (\textit{Cumulative Distribution Function }- CDF) de uma variável aleatória de valor real $x$ é a função dada por: 
\[
F_{X}\left(x\right)=P\left(X\leq x\right)
\]
Ou seja, temos a probabilidade de uma variável aleatória $X$ assumir um valor menor ou igual a $x$. A CDF de uma variável contínua $X$ pode ser expressa em termos de uma função densidade $f_{x}$ de probabilidade dada por: 
\[
F_{X}\left(x\right)=\int^{x}_{-\infty}f_{X}\left(t\right)dt
\]

A CCDF pode ser definida, então, como: 
\[
\overline{F}_{X}\left(x\right)=P\left(>x\right)=1-F_{X}\left(x\right)
\]
Isto é, a probabilidade de uma variável aleatória $X$ assumir um valor maior que $x$. No nosso caso, estamos interessados na probabilidade de que, ao selecionarmos um agente aleatoriamente, ele tenha uma riqueza superior a $w$, fazendo uma pequena mudança de notação, temos:

\[
\overline{F}_{X}\left(w\right)=P\left(X>w\right)=1-F_{X}\left(w\right)
\]

Podemos calcular ,então, a CDF de Pareto da seguinte forma:

\[
F_{X}\left(w\right)=Pr\left(X\leq w\right)=\int^{w}_{-\infty}p\left(w'\right)dw'
\]

Resolvendo a integral: 
\[
\begin{split}F_{X}\left(w\right) & =\int^{w}_{\beta}\frac{\alpha\beta^{\alpha}}{x'^{\alpha+1}}dw'\\
 & =\alpha\beta^{\alpha}\left(\frac{w'^{-\left(\alpha+1\right)+1}}{-\left(\alpha+1\right)+1}\right)|^{w}_{\beta}\\
 & =\alpha\beta^{\alpha}\left(\frac{w'^{-\alpha}}{-\alpha}\right)|^{w}_{\beta}\\
 & =\alpha\beta^{\alpha}\left(\frac{w^{-\alpha}}{-\alpha}-\frac{\beta^{-\alpha}}{-\alpha}\right)\\
 & =\left(\frac{\beta}{\beta}\right)^{\alpha}-\left(\frac{\beta}{w}\right)^{\alpha}\\
 & =1-\left(\frac{\beta}{w}\right)^{\alpha}
\end{split}
\]
Então, o CCDF da lei de Pareto será:

\[
\overline{F}_{X}\left(w\right)=1-\left(1-\left(\frac{\beta}{w}\right)^{\alpha}\right)=\left(\frac{\beta}{w}\right)^{\alpha}=aw^{-\alpha}
\]

Onde $a=\beta^{\alpha}$. Se aplicarmos o logaritmo natural, temos:

\[
\ln\overline{F}_{X}\left(w\right)=\ln\left(aw^{-\alpha}\right)=\ln\left(a\right)+\ln\left(w^{-\alpha}\right)=\ln\left(a\right)-\alpha\ln\left(w\right)
\]
Onde retomamos a equação de Preto $\ln N\left(x\right)=\ln A-\tilde{\gamma}\ln\left(x\right)$, identificando $\overline{F}(w)=N(x)$ e $\alpha=\tilde{\gamma}$. Porém, estamos interessados em obter o índice de Pareto de uma distribuição de valores. Isso é obtido mais facilmente após realizarmos uma linearização. Para linearizar o CCDF da lei de Pareto, começamos aplicando o logaritmo:

\[
\begin{split}\log_{10}\left(\overline{F}_{X}\left(w\right)\right) & =\log_{10}\left(aw^{-\alpha}\right)\\
 & =\log_{10}\left(a\right)+\log_{10}\left(w^{-\alpha}\right)
\end{split}
\]
Como $a$ é um valor constante, Nada nos impede de escrever $a=10^{b}$ . Escrevendo a linearização do CCDF de Pareto como $y=\log_{10}\left(\overline{F}_{X}\left(w\right)\right)$, temos então:

\[
\begin{split}y & =\log_{10}\left(10^{b}\right)+\log_{10}\left(w^{-\alpha}\right)\\
 & =b\log_{10}\left(10\right)-\alpha\log_{10}\left(w\right)
\end{split}
\]
Fazendo uma substituição de variável do tipo $x=\log_{10}\left(w\right)$ , ficamos então com a seguinte equação linear:

\[
y=b-\alpha x
\]
O índice de Pareto nada mais é do que a inclinação desta reta. Além disso, podemos escrever a constante $b$ em termos das constantes originais da lei de Pareto da seguinte forma:

\[
b=\log_{10}\left(a\right)=\log_{10}\left(\beta^{\alpha}\right)
\]
Ou ainda, de forma inversa, podemos obter $\beta$ a partir da constante $b$ e do coeficiente $\alpha$ calculando $\beta=10^{\left(b/\alpha\right)}$. Este é o CCDF linearizado da lei de Pareto , após aplicarmos o logaritmo na base 10, isto precisará ser retomado mais tarde. Para calcular o índice de Pareto de um conjunto de dados é comum aplicar algum método matemático a esse conjunto de dados, tipicamente regressão linear.

\subsection{Gini}

A curva de Lorenz é dada por uma função do tipo $L(x)$ que nos indica a fração de riqueza (ou renda) em posse da fração $x$ da população. A igualdade perfeita é dada, então, por uma reta $L(x)=x$.

O coeficiente de Gini é uma medida da desigualdade em um conjunto de dados e é dado pela razão entre a área que se encontra entre a reta de igualdade perfeita e a curva de Lorenz, e a área total sob a reta de igualdade perfeita.

Se temos um conjunto de dados ordenados $(y_{1}\leq y_{2}\leq\dots\leq y_{N})$, a proporção acumulada da população, ou seja, a posição no eixo $X$ para o dado $i$ é dada por:

\[
x_{i}=\frac{i}{N}
\]

Consequentemente a distância entre dois dados $i$ e $i+1$ é dada por

\[
\Delta x=\frac{1}{N}
\]

A altura da curva de Lorentz é dado por:

\[
y_{i}=\frac{\sum^{i}_{k=1}y_{k}}{\sum^{N}_{k=1}y_{k}}
\]

Para a reta de igualdade perfeita temos: 
\[
y^{*}_{i}=x_{i}=\frac{i}{N}
\]

Logo a diferença entre as alturas é dada por:

\[
\Delta y_{i}=\frac{i}{N}-\frac{\sum^{i}_{k=1}y_{k}}{\sum^{N}_{k=1}y_{k}}
\]

Assim sendo, a área dada pela diferença entre a reta da igualdade e a curva de Lorenz para a distância entre dois dados é dada por:

\[
\Delta x\Delta y_{i}=\frac{1}{N}\left(\frac{i}{N}-\frac{\sum^{i}_{k=1}y_{k}}{\sum^{N}_{k=1}y_{k}}\right)
\]

Somando então sobre todos os dados:

\[
\begin{split}A & =\sum^{N}_{i=1}\frac{1}{N}\left(\frac{i}{N}-\frac{\sum^{i}_{k=1}y_{k}}{\sum^{N}_{k=1}y_{k}}\right)\\
 & =\frac{1}{N}\left(\frac{1}{N}\sum^{N}_{i=1}i-\frac{\sum^{N}_{i=1}\sum^{i}_{k=1}y_{k}}{\sum^{N}_{k=1}y_{k}}\right)
\end{split}
\]
O primeiro termo entre parênteses: 
\[
\frac{1}{N}\sum^{N}_{i=1}i=\frac{1}{N}\frac{N\left(N+1\right)}{2}=\frac{N+1}{2}
\]
Para o somatório duplo no segundo termo, percebendo que para cada $k$ o valor $y_{k}$ vao aparecer em todos $i\geq k$, isto é $N-k+1$ vezes: 
\[
\sum^{N}_{i=1}\sum^{i}_{k=1}y_{k}=\sum^{N}_{k=1}\left(N-k+1\right)y_{k}
\]
Então: 
\[
A=\frac{1}{N}\left(\frac{N+1}{2}-\frac{\sum^{N}_{k=1}\left(N-k+1\right)y_{k}}{\sum^{N}_{k=1}y_{k}}\right)
\]
Por fim, nos resta apenas lembrar que o Gini é dado pela razão entre a área que se encontra entre a reta da igualdade e a curva de Lorenz ($A$) e a própria área abaixo da reta da igualdade. Considerando que os dados estejam normalizados, a área abaixo da reta da igualdade perfeita corresponde simplesmente à metade da área de um quadrado de lado $1$ ou seja, $1/2$ o coeficiente de Gini é então:

\[
G=\frac{2}{N}\left(\frac{N+1}{2}-\frac{\sum^{N}_{k=1}\left(N-k+1\right)y_{k}}{\sum^{N}_{k=1}y_{k}}\right)
\]
Ou em uma forma mais comumente encontrada na literatura\cite{gini}: 
\[
G=\frac{1}{N}\left(N+1-2\left(\frac{\sum^{N}_{k=1}\left(N-k+1\right)y_{k}}{\sum^{N}_{k=1}y_{k}}\right)\right)
\]

\subsection{Revisão de conceitos estatísticos}

Vamos fazer uma revisão baseada no apêndice do livro de Econometri Básica\cite{gujarati2011econometria}.

\textbf{Somatório:}
\begin{itemize}
\item $\sum^{n}_{i=1}ck=c\sum^{n}_{i=1}k=ckn$
\end{itemize}
\textbf{Espaço amostral:}
\begin{itemize}
\item \textbf{Espaço amostral (população)}: O conjunto de todos os resultados possíveis de um experimento aleatório.
\begin{itemize}
\item \textbf{Ponto amostral}: cada membro desse espaço amostral.
\end{itemize}
\item \textbf{Evento}: um subconjunto do espaço amostral. 
\begin{itemize}
\item \textbf{Eventos mutuamente exclusivos}: se a ocorrência de um evento eliminar a ocorrência do outro. 
\end{itemize}
\item \textbf{Eventos (coletivamente) exaustivos}: se o conjunto de eventos exaurem todas os possíveis resultados de um experimento. 
\end{itemize}
\textbf{Probabilidade e variáveis aleatórias}

Seja $A$ um evento em um espaço amostral, por $P(A)$, a probabilidade do evento $A$, entendemos a proporção de vezes que o evento A ocorrerá em repetidas tentativas de um experimento. Ou, em um total de $n$ possíveis resultados igualmente prováveis de um experimento, se $m$ deles são favoráveis à ocorrência do evento A, definimos a razão $m/n$ como a frequência relativa de A. 
\begin{itemize}
\item $0\leq P\left(A\right)\leq1$ para cada $A$
\item Se $A,B,C$ são um conjunto de eventos exaustivos então a probabilidade de termos um evento A ou B ou C ... é $P\left(A+B+C\right)=1$.
\item Se A e B são mutuamente exclusivos: $P\left(A+B\right)=P\left(A\right)+P\left(B\right)$
\item \textbf{Variável aleatória}: uma variável cujo valor é determinado pelo resultado de um experimento aleatório, usualmente denotado por letras maíusculas, ex: $X$.
\begin{itemize}
\item Os valores assumidos pela variável aleatória são normalmente denotados por letras minúsculas, ex.: $x$.
\end{itemize}
\end{itemize}
\textbf{Função de densidade de probabilidade}

\textbf{Discreto}: Se $X$ é uma variável aleatória discreta $f\left(x\right)=P\left(X=x_{i}\right)$ então a função densidade de probabilidade (FDP) significa a probabilidade de que a variável aleatória discreta $X$ tome o valor $x_{i}$, sendo $f\left(x\right)=0$ para $x\neq x_{i}$.

\textbf{Contínuo}: o elemento da probabilidade $f\left(x\right)dx$ é a probabilidade associada a um pequeno intervalo de uma variável contínua e $P\left(a\leq x\leq b\right)$ indica a probabilidade de que $x$ situe-se no intervalo entre $a$ e $b$.

Sendo:

\begin{align*}
f\left(x\right) & \geq0\\
\int^{+\infty}_{-\infty}f\left(x\right)dx & =1\\
\int^{b}_{a}f\left(x\right)dx & =P\left(a\leq x\leq b\right)
\end{align*}

\textbf{Função de densidade de probabilidade conjunta: }

\textbf{Caso discreto} : É a probabilidade de que $X$ tome valor $x_{i}$ e $Y$ tome o valor $y_{i}$ dada por $f\left(x,y\right)=P\left(X=x_{i}\text{ e }Y=y_{i}\right)$.

\textbf{Caso contínuo:}

\begin{align*}
f\left(x,y\right) & \geq0\\
\int^{+\infty}_{-\infty}\int^{+\infty}_{-\infty}f\left(x,y\right)dxdy & =1\\
\int^{b}_{a}\int^{d}_{c}f\left(x,y\right)dxdy & =P\left(a\leq x\leq b,c\leq y\leq d\right)
\end{align*}

\textbf{Função de densidade de probabilidade marginal (ou individual)}

\textbf{Caso discreto}: em relação a $f\left(x,y\right)$, $f\left(x\right)$ e $f\left(y\right)$ são obtidas por $f\left(x\right)=\sum_{y}f\left(x,y\right)$, onde somamos sobre todos valores de $Y$.

\textbf{Função de densidade de probabilidade condicional:}

Expressa a probabilidade de que $X$ assuma o valor $x$, posto que $Y$ assumiu o valor $y$:

\[
f\left(x|y\right)=P\left(X=x|Y=y\right)
\]

Se relaciona com a função de densidade de probabilidde conjunta da seguinte forma:

\[
f\left(x|y\right)=\frac{f\left(x,y\right)}{f\left(y\right)}\rightarrow f\left(x,y\right)=f\left(y\right)f\left(x|y\right)
\]

A probabilidade de que $x$ e $y$ aconteçam é o produto entre a probabilidade que $y$ aconteça e a probabilidade de que $x$ aconteça posto que $y$ aconteceu.

\textbf{Independência estatística:}

As duas variáveis são independente ssomente se:

\[
f\left(x,y\right)=f\left(x\right)f\left(y\right)
\]
\textbf{Média (valor esperado) de uma dsistribuição de probabilidde:}

Propriedades:
\begin{itemize}
\item Se $c$ é constante $E\left(b\right)=b$.
\item Se $k$ e $c$ são constantes: $E\left(aX+b\right)=aE\left(x\right)+b$
\item Se $X$ e $Y$ são variáveis aleatórias independentes $E\left(XY\right)=E\left(X\right)E\left(Y\right)$\footnote{Isso não implica que E$\left(\frac{X}{Y}\right)=\frac{E\left(X\right)}{E\left(Y\right)}$.}
\item Se $X$ é uma variável aleatória com função função de densidade de probabilidade $f\left(x\right)$ e $g\left(x\right)$ é qualquer função de $X$ então:
\end{itemize}
\begin{align*}
E\left[g\left(X\right)\right] & =\sum_{x}g\left(x\right)f\left(x\right) & \text{se {X} for discreta}\\
 & =\int^{+\infty}_{-\infty}g\left(x\right)f\left(x\right)dx & \text{se {X} for contínua}
\end{align*}

Note que se $g\left(x\right)=x$ então temos simplesmente o valor esperado de $X$.

\textbf{Variância}: É a distribuição (dispersão) em torno do valor esperado.

\[
var\left(x\right)=\sigma^{2}_{x}=E\left[\left(X-\mu\right)^{2}\right]
\]

Onde $E\left(X\right)=\mu$ é o valor esperado, e $\sigma_{x}$também é conhecido como desvio padrão.

\begin{align*}
var\left(X\right) & =\sum_{x}\left(x-\mu\right)^{2}f\left(x\right) & \text{se {X} for discreta}\\
 & =\int^{+\infty}_{-\infty}\left(x-\mu\right)^{2}f\left(x\right)dx & \text{se {X} for contínua}
\end{align*}

Propriedades:
\begin{itemize}
\item $var=E\left(X^{2}\right)-E\left(X\right)^{2}$
\item Se $c$ é constante a variância é $0$
\item Se $a$ também é constante $var\left(aX+c\right)=a^{2}var\left(X\right)$
\item Se $X$ e $Y$ são variáveis aleatórias independentes $var\left(X\pm Y\right)=var\left(X\right)+var\left(Y\right)$
\end{itemize}
\textbf{Covariância}

A covariância de duas variáveis aleatória $X$ e $Y$ com média $\mu_{x}$ e $\mu_{y}$ é dado por:

\[
cov\left(X,Y\right)=E\left[\left(X-\mu_{x}\right)\left(Y-\mu_{y}\right)\right]
\]

Ou seja:

\begin{align*}
cov\left(X,Y\right) & =\sum_{x,y}\left(x-\mu_{x}\right)\left(y-\mu_{y}\right)f\left(x,y\right)\\
 & =\sum_{x,y}xyf\left(x,y\right)-\mu_{y}\sum_{x,y}xf\left(x,y\right)-\mu_{y}\sum_{x,y}yf\left(x,y\right)+\mu_{y}\mu_{x}\sum_{x,y}f\left(x,y\right)\\
 & =\sum_{x,y}xyf\left(x,y\right)-\mu_{y}\sum_{x}xf\left(x\right)-\mu_{y}\sum_{y}yf\left(y\right)+\mu_{y}\mu_{x}\\
 & =\sum_{x,y}xyf\left(x,y\right)-\mu_{y}\mu_{x}-\mu_{x}\mu_{y}+\mu_{y}\\
 & =\sum_{x,y}xyf\left(x,y\right)-\mu_{y}\mu_{x}
\end{align*}

Ou de forma análoga pro caso contínuo:

\begin{align*}
cov\left(X,Y\right) & =\int^{+\infty}_{-\infty}\left(x-\mu_{x}\right)\left(y-\mu_{y}\right)f\left(x,y\right)dxdy\\
 & =\int^{+\infty}_{-\infty}xyf\left(x,y\right)dxdy-\mu_{y}\mu_{x}
\end{align*}

De forma geral:

\[
cov\left(X,Y\right)=E\left(XY\right)-\mu_{y}\mu_{x}
\]

A covariância mede o quanto as duas variáveis variam conjuntamente em relação às médias. Se $a=\left(x-\mu_{x}\right)>0$ e $b=\left(y-\mu_{y}\right)>0$ então $ab>0$. Agora se ambos $a<0$ também $b<0$ então também temos $ab>0$. Ou seja, se sempre variam juntas vamos ter predominância de $ab>0$ resultando em uma covariância positiva, de forma similar se sempre variam em sentido oposto, teremos predominâmcia de $ab<0$ então teremos uma covariância negativa. Mas se for igualmente distribuído, então teremos uma covariância próxima de zero.

Propriedades:
\begin{itemize}
\item Se $X$ e $Y$ são independentes então $cov\left(X,Y\right)=0$.
\item Se $a,b,c,d$ são constantes: $cov\left(a+bX,c+dY\right)=bd\cdot cov\left(X,Y\right)$
\end{itemize}
\textbf{Coeficiente de correlação:}

\[
\rho=\frac{cov\left(X,Y\right)}{\sqrt{var\left(X\right)var\left(Y\right)}}=\frac{cov\left(X,Y\right)}{\sigma_{x}\sigma_{y}}
\]

\textbf{Outra interpretação}: Vamos olhar para a correlação de Pearson para uma amonstra constituída por um conjunto de pontos $\left\{ \left(x_{i},y_{i}\right)\right\} $, acredito que a melhor forma é ter uma interpretção geométrica. Se separarmos em um conjunto de pontos $X=\left(x_{1},x_{2},\dots,x_{n}\right)$ e outro $Y=\left(y_{1},y_{2},\dots,y_{n}\right)$ então se descontarmos os valores médios:

\[
\tilde{X}=\left(x_{1}-\overline{x},x_{2}-\overline{x},\dots,x_{n}-\overline{x}\right),\quad,\tilde{Y}=\left(y_{1}-\overline{y},y_{2}-\overline{y},\dots,y_{n}-\overline{y}\right)
\]
Vamos sobrepor a distribuição de pontos em torno da origem. Agor podemos analisar como cada uma das distribuições se afastam da origem. Se teos só três pontos, temos que $\tilde{X}$ e $\tilde{Y}$ denotam vetores em $3D$. O produto escalar de ambos é dado por:

\[
\tilde{Y}\cdot\tilde{X}=\left\Vert \tilde{X}\right\Vert \left\Vert \tilde{Y}\right\Vert \cos\left(\theta\right)
\]

Então:

\[
\cos\left(\theta\right)=\frac{\tilde{Y}\cdot\tilde{X}}{\left\Vert \tilde{X}\right\Vert \left\Vert \tilde{Y}\right\Vert }=\frac{\sum^{n}_{i=1}\left(x_{i}-\overline{x}\right)\left(y_{i}-\overline{y}\right)}{\sqrt{\sum^{n}_{i=1}\left(x_{i}-\overline{x}\right)^{2}}\sqrt{\sum^{n}_{i=1}\left(y_{i}-\overline{y}\right)^{2}}}=r
\]

Portanto, a correlação de Pearson nos fornece o cosseno do ângulo entre os vetores formados pelas variáveis após removermos seus valores médios. Em outras palavras, ela mede o alinhamento entre as variações de $X$ e $Y$ em torno de suas respectivas médias, interpretadas como vetores em um espaço n-dimensional. Conectando com a discussão anterior, tínhamos:

\[
\rho=\frac{cov(X,Y)}{\sigma_{x}\sigma_{y}}=\frac{\sum_{x,y}(x-\mu_{x})(y-\mu_{y})f(x,y)}{\sqrt{\sum_{x}(x-\mu_{x})^{2}f(x)}\sqrt{\sum_{y}(y-\mu_{y})^{2}f(y)}}
\]

Neste caso, como temos $n$ pares observados, a probabilidade empírica de selecionar qualquer observação individual de $X$ (ou de${Y})$ é $f(x_{i})=f(y_{i})=\frac{1}{n}.$ Além disso, como os dados são observados em pares $(x_{i},y_{i})$, a distribuição conjunta empírica satisfaz: $f(x_{i},y_{j})=\begin{cases}
\frac{1}{n}, & i=j\\
0, & i\neq j
\end{cases}$. Assim, o termo da covariância pode ser escrito como:

\[
\sum_{x,y}(x-\mu_{x})(y-\mu_{y})f(x,y)=\frac{1}{n}\sum^{n}_{i=1}(x_{i}-\mu_{x})(y_{i}-\mu_{y}).
\]

Os fatores $1/n$ do numerador e do denominador então se cancelam, resultando em:

\[
\rho=\frac{\sum^{n}_{i=1}(x_{i}-\bar{x})(y_{i}-\bar{y})}{\sqrt{\sum^{n}_{i=1}(x_{i}-\bar{x})^{2}}\sqrt{\sum^{n}_{i=1}(y_{i}-\bar{y})^{2}}}
\]

\textbf{Spearman}

Não difere muito do caso anterior, apens que ao invés de trabalharmos com os dados originais, nós transformarmos eles em um ranking, só então aplicamos a correlação de Pearson.

Tanto nesta seção quanto na anterior, devemos ter em mente que não analisamos os erros padrão nem os p-values. Em geral, essa análise adicional é importante para avaliar a significância estatística dos resultados, especialmente nos testes de correlação. No entanto, em muitos casos a regressão linear pode ser utilizada apenas como uma ferramenta descritiva, isto é, com o objetivo de encontrar a reta que melhor representa a distribuição dos pontos, independentemente de inferência estatística sobre os coeficientes.

\textbf{Variâncias de variáveis correlacionadas}

Se são correlacionadas então: $var\left(X\pm Y\right)=var\left(X\right)+var\left(Y\right)\pm cov\left(X,Y\right)$

\textbf{Expectativa condicional}

A expectativa condicional de $X$ dado $Y=y$ é:

\begin{align*}
E\left(X|Y=y\right) & =\sum_{x}xf\left(x|Y=y\right) & \text{se {X} for discreta}\\
 & =\int^{+\infty}_{-\infty}xf\left(x|Y=y\right)dx & \text{se {X} for contínua}
\end{align*}

Se $E\left(X|Y\right)$ é uma variável aleatória porque é uma função da variável condicionamente, $E\left(X|Y=y\right)$ é uma constante para um dado valor específico de $Y$.

\textbf{Variância condicional}

De modo análogo sendo $var\left(X|Y=y\right)=E\left\{ \left[X-E\left(X|Y=y\right)\right]^{2}|Y=y\right\} $:

\begin{align*}
E\left(X|Y=y\right) & =\sum_{x}\left[X-E\left(X|Y=y\right)\right]^{2}f\left(x|Y=y\right) & \text{se {X} for discreta}\\
 & =\int^{+\infty}_{-\infty}\left[X-E\left(X|Y=y\right)\right]^{2}f\left(x|Y=y\right)dx & \text{se {X} for contínua}
\end{align*}

Propriedades da expectativa e variância condicional:
\begin{itemize}
\item Se $f\left(X\right)$ então $E\left(f\left(X\right)|X\right)=f\left(X\right)$
\item Se $f\left(X\right)$ e $g\left(X\right)$ são funções de $X$ então $E\left(f\left(X\right)Y+g\left(X\right)\right)=f\left(X\right)E\left(Y|X\right)+g\left(X\right)$
\item Lei das expectativas iteradas: $E\left(Y\right)=E_{X}\left[E\left(Y|X\right)\right]=\sum_{x}E\left(Y|X=x\right)f\left(x\right)$, onde $E_{X}$ denota que a expectativa ta cobrindo todos os valores de $X$.
\item Se $X$ e $Y$ forem independentes $E\left(Y|X\right)=E\left(Y\right)$ e $var\left(Y|X\right)=var\left(Y\right)$
\end{itemize}
\textbf{Momentos de ordem superior:}

Momento de ordem $r$: $E\left[\left(X-\mu\right)^{r}\right]=m_{r}$. Terceiro e quarto momento são utilizados no estudo da ``forma'' de uma probabilidade. Para o calculo da assimetria (falta de simetria) calculamos $S=m_{3}/\sigma^{3}$ e curtose (elevação ou achatamento) $K=m_{4}/m^{2}_{2}$. Nese caso $K<3$ são chamadas platiúrticas (gordas ou de cada curtas) e $K>3$são chamadas de leptocúrticas (magras ou cadas longas). E ainda $K=3$ é chamada de mesocúritca (distribuição normal).

\textbf{Distribuições:}

Sobre a normal:
\begin{itemize}
\item Por questão de notação, denotamos uma variável distribuída de forma normal como$X\sim N\left(\mu,\sigma^{2}\right)$.
\item Uma combinação linear de variáveis de distribuição normal é distribuída normalmente: $Y=aX_{1}+bX_{2}\sim N\left[\left(a\mu_{1}+b\mu_{2}\right),\left(a^{2}\sigma^{2}_{1}+b^{2}\sigma^{2}_{2}\right)\right]$
\item Teorema central do lime: $\overline{X}_{x\rightarrow\infty}\sim N\left(\mu,\frac{\sigma^{2}}{n}\right)$, isto é, $\overline{X}$ aproxima-se da distribuição normal com média $\mu$ e variância $\sigma^{2}/n$. 
\begin{itemize}
\item Esse resultado é verdadeiro não importando a forma da FDP original.
\end{itemize}
\item Se $X$ e $Y$ possuem distribuição conjunta normal (isto é, os pares $(X,Y)$ formam uma distribuição normal bidimensional), então toda a dependência entre elas é determinada pela covariância. Nesse caso, dizer que:$Cov(X,Y)=0$ é equivalente a dizer que são independentes.
\end{itemize}
\textbf{Inferência estatística: estimação}

Pode ser que sabemos (ou supomos) que uma variável aleatória $X$ segue uma distribuição de probabilidade particicular, mas não sabemos os valores dos parâmetros da distribuição. Por exemplo, se temos uma distribuição normal, podemos querer saber a média e a variância.O procedimento habitual é ter uma amostra aleatória de tamanho $n$ e então utilizar estes dados para estimar os parâmetros desconhecidos. Se temos uma variável aleatória $X$ com FDP $f(x)$ e repetimos o experimento 100 vezes de forma independente, podemos considerar este como um único experimento que tem como resultado uma amostra $X_{1},X_{2},\dots,X_{100}$, onde todas tem a mesma distribuição $f(x)$. Nesse caso, a FDP conjunta deste experimento é $f(x_{1},\dots,x_{100})=f(x_{1})\cdots f(x_{100}),$pois as observações são independentes e identicamente distribuídas. Assim, $X_{1},\dots,X_{100}$ constituem uma amostra aleatória de tamanho $n=100$ proveniente de uma população com FDP $f(x)$. Deve estar claro: o experimento aqui consiste em retirar 100 valores aleatórios $x_{i}$de $X$. Ou seja, uma amostra aleatória de tamanho $n$ é o conjunto de $n$ valores observados de uma variável aleatória gerada repetidamente segundo a mesma FDP $f(x)$.

\textbf{Estimação pontual}

Vamos supor que sabemos a FDP, mas não sabemos um parâmetro $\theta$. Sorteamos então uma amostra aleatória de tamanho $n$ e desenvolvemos uma função dos valores da amostra de modo que $\widehat{\theta}=f\left(x_{1},\dots,x_{n}\right)$ (chmado de estimador ou estatística) forneça uma estimativa do verdadeiro $\theta$ (chamado de estimativa)\footnote{Vale dizer que $\widehat{\theta}$ também pode ser tratada como uma variável aleatória pois é uma função dos dados amostrais, com uma própria distribuição.}. Se ele fornece uma estimativa única, é então conhecido como estimativa pontual.

\textbf{Estimação intervalar}

Agora não temos mais apenas um valor estimado, mas um intervalo no qual esperamos que o verdadeiro valor do parâmetro esteja. O conceito importante é o da distribuição de probabilidade do estimador. Por exemplo, se uma variável aleatória $X$ possui distribuição normal, então a média da amostra também possui uma distribuição normal com média $\mu$ (a média verdadeira da população) e variância $\sigma^{2}/n$. Assim, a distribuição do estimador $\overline{X}$ é $\overline{X}\sim N\left(\mu,\frac{\sigma^{2}}{n}\right).$

Como sabemos que, em uma distribuição normal, o intervalo entre $\mu\pm2\sigma/\sqrt{n}$ corresponde a aproximadamente $95\%$ da área sob a curva da distribuição, temos que $P\left(\mu-\frac{2\sigma}{\sqrt{n}}\leq\overline{X}\leq\mu+\frac{2\sigma}{\sqrt{n}}\right)\approx0.95,$sendo $\mu$ a média da distribuição da população e também da distribuição teórica do estimador, e $\overline{X}$ a média empírica da amostra. Ou seja, temos aproximadamente $95\%$ de chance de que nossa média amostral $\overline{X}$ esteja a uma distância inferior a $\frac{2\sigma}{\sqrt{n}}$ da média verdadeira teórica. Podemos mostrar facilmente que $a-b\leq X\leq a+b$ é equivalente a $X-b\leq a\leq X+b$, bastando rearranjar os termos da desigualdade, portanto, $P\left(\overline{X}-\frac{2\sigma}{\sqrt{n}}\leq\mu\leq\overline{X}+\frac{2\sigma}{\sqrt{n}}\right)\approx0.95$\footnote{Mas vale destacar que aqui $\sigma$ representa o desvio padrão verdadeiro da população. Na prática, como esse valor geralmente é desconhecido, podemos aproximá-lo pelo desvio padrão medido empiricamente na amostra. Contudo, essa substituição adiciona uma incerteza extra à estimação, já que o próprio desvio padrão passa a ser estimado a partir da amostra.}. Equivalentemente, o intervalo $\left[\overline{X}-\frac{2\sigma}{\sqrt{n}},\overline{X}+\frac{2\sigma}{\sqrt{n}}\right]$ possui 95\% de chances de conter o verdadeiro valor de $\mu$\footnote{Detalhe técnico: Em mensurações repetidas como essa, os intervalos assim estabelecidos incluirão o verdadeiro $\mu$ com 95\% de confiança. Porém não podemos dizer que seja de 95\% a probabilidade de que um intervalo particular medido contenha $\mu$. Como esse intervalo é fixado, a probabilidade de que ele incluirá $\mu$ é 0 ou 1. O que podemos afirmar é que, se construirmos 100 desses intervalos, 95 dos 100 intervalos incluirão $\mu$; não podemos garantir que um intervalo em particular incluirá necessarimante $\mu$.}.

De forma mais geral, na estimação intervalar, construímos dois estimadores $\widehat{\theta}_{1}$ e $\widehat{\theta}_{2}$ de maneira que:

\[
P\left(\widehat{\theta}_{1}\leq\theta\leq\widehat{\theta}_{2}\right)=1-\alpha\quad0<\alpha<1
\]

Ou seja, podemos afirmar tanto que $p=1-\alpha$ é a probabilidade de que este intervalo contenha o verdadeiro $\theta$ . Este intervalo se chama intervalo de confiança $p$, sendo $p$ o coeficiente de confiança e $\alpha$ o nível de significância. 

\textbf{Métodos de estimação}

Há três métodos principais de estimação de parâmetros.

\textbf{Método dos mínimos quadrados:} consiste em definir uma função erro, geralmente baseada na soma dos quadrados dos desvios entre os dados observados e os valores previstos pelo modelo, e minimizá-la. Já utilizamos essa ideia anteriormente para estimar a inclinação de uma reta. 

\textbf{Método da verossimilhança}: depende da função de verossimilhança. Vamos supor que temos uma variável aleatória $X$ com uma FDP $f\left(X,\theta\right)$, novamente conhecemos a FDP nãos não o valor do parâmetro. Se temos uma amostra de $n$ variáveis aleatórias $X_{i}$, novamente temos uma FDP conjunta. Agora temos uma interpretação dual desta FPD.

Se sabemos o $\theta$ ela é uma FDP de observar os dados dos valores amostrais. Mas mas se temos os valores amostrais $x_{n}$ então ela pode ser uma função de $\theta$. Neste caso chamamos de função de verossimilhança, e agora buscamos o valor do parâmetro $\theta$ que torna a amostra observada a mais plausível possível. A função verossimilhança, é então:

\[
L\left(\theta;x_{1},\dots,x_{n}\right)=\prod_{i}f\left(x_{i};\theta\right)
\]

O estimador de máxima verossimilhança de $\theta$ é aquele que maximiza a função de verossimilhança da amostra. Por questões matematicas em geral tomamos o logaritmo, e seguindo as regras de cálculo diferenciamos a função com respeito a incognita e igualamos a derivada a zero. O estimador obtido assim é chamado de estimador de máxima verossimilhança. Por exemplo, a distribuição de Poisson é $f\left(X\right)=e^{-\lambda}\lambda^{x}/x!$ e $E\left(X\right)=\lambda$. Então obtemos $n$ valores a partir da FDP e construímos uma amostra, a FDP conjunta da amostra é:

\[
L\left(\theta;x_{1},\dots,x_{n}\right)=\prod_{i}\frac{e^{-\lambda}\lambda^{x_{i}}}{x_{i}!}
\]

Tirando o logaritmo:

\begin{align*}
\log L & =\sum\left[-\lambda\log e+x_{i}\log\lambda-\log x_{i}!\right]\\
 & =-\lambda n+\sum x_{i}\log\lambda-\sum\log x_{i}!
\end{align*}

Diferenciando com respeito a $\lambda$ e igualando a zero:

\[
0=-n+\frac{\sum x_{i}}{\widehat{\lambda}}
\]

Então $\widehat{\lambda}=\frac{\sum x_{i}}{n}=\overline{X}$, sendo que $\overline{X}$ é o estimador de $E\left(X\right)$.

\textbf{Método dos momentos: }O método dos momentos baseia-se no chamado princípio da analogia, segundo o qual as propriedades observadas na amostra devem reproduzir, ao menos aproximadamente, as propriedades teóricas da população. Em termos matemáticos, isso significa substituir os momentos (como a média $E(X)$ por exemplo) pelos seus correspondentes amostrais (como a média $\overline{X}$). Assim, os parâmetros desconhecidos da distribuição são estimados impondo que os momentos teóricos coincidam com os momentos observados na amostra. A justificativa desse procedimento decorre da Lei dos Grandes Números, segundo a qual a média amostral converge para a média populacional à medida que o tamanho da amostra aumenta.

Por exemplo, considere uma variável aleatória com distribuição exponencial, $X\sim\text{Exp}(\lambda)$. Nesse caso, a média populacional é dada por $E(X)=\frac{1}{\lambda}$. Pelo método dos momentos, substituímos a média populacional pela média observada na amostra, isto é, $\overline{X}=\frac{1}{\lambda}$ e obtemos o estimador pelo método dos momentos: $\hat{\lambda}=\frac{1}{\overline{X}}$. Assim, o parâmetro desconhecido da distribuição é estimado a partir da média observada na amostra. Mas eu não me aprofundarei nem no princípio da analogia, e nem no método dos momentos generalizados.

\textbf{Propriedades de pequenas amostras:} Por trás desses conjuntos de propriedades está a noção de que um estimador possui uma distribuição de probabilidade.

\textbf{Sem viés}: $E\left(\widehat{\theta}\right)=\theta$, sendo $\text{viés}=E\left(\widehat{\theta}\right)-\theta$

\textbf{Variância mínima}: dizemos que $\widehat{\theta}$ é um estimador de variância mínima se sua variância for menor que a de qualquer outro estimador.

\textbf{Melhor estimador não enviesado} \textbf{(ou estimador eficiente)}: Se é o estimador de menor variância e sem viés.

\textbf{Linearidade}: Um estimador $\widehat{\theta}$ é um estimador linear se ele é um função linear das observações da amostra. Como por exemplo, a média da amostra é uma função linear dos valores de $X$ :$\overline{X}=\frac{1}{n}\sum X_{i}$.

\textbf{Melhor estimador linear não viesado}: se o estimador é linear, não viesado e possui uma variância mínima.

\textbf{Estimador com erro quadrado médio mínimo (MSE)}: Definimos $MSE=E\left[\left(\widehat{\theta}-\theta\right)\right]^{2}$

Enquanto a variância mensura a dispersão da distribuição de $\widehat{\theta}$ em torno da média, o MSE mensura a dispersão em torno do valor verdadeiro. E relação pode ser escrita como:

\[
MSE=E\left[\left(\widehat{\theta}-\theta\right)^{2}\right]=var\left(\widehat{\theta}\right)+\left[vies\left(\widehat{\theta}\right)\right]^{2}
\]

Na prática, o critério de mínimo erro quadrático médio (MSE) é frequentemente utilizado quando se deseja equilibrar simultaneamente viés e variância, permitindo inclusive a utilização de estimadores viesados caso isso resulte em menor erro total de estimação.

\textbf{Propriedades de grandes amostras ou assintóticas:}

Um estimador pode não satisfazer uma ou mais propriedades em amostra pequena, mas à medida que o tamanho da amostra cresce, o estimador pode possuir várias propriedade estatísticas desejáveis. 

\textbf{Ausência assintótica de viés}: $\lim_{n\rightarrow\infty}E\left(\widehat{\theta}_{n}\right)=\theta$

\textbf{Consistência}: Quando o estimador $\widehat{\theta}$ se aproxima do valor verddeiro de $\theta$ à medida que o tamanho da amostra torna-se cada vez maior. Diz-se, mais formalmente, que $\widehat{\theta}$ é um estimador consistente de $\theta$ se a probabilidade de que o valor absoluto da diferença entre $\widehat{\theta}$ e $\theta$ seja menor do que $\delta$, isto é, uma quantidade positiva arbitrariamente pequena, aproxima-se de 1 quando $n$ tende ao infinito. Simbolicamente:

\[
\lim_{n\rightarrow\infty}P\left\{ \left|\widehat{\theta}-\theta\right|<\delta\right\} \quad\delta>0
\]

Isso é frequentemente expresso como:

\[
\underset{n\rightarrow\infty}{\text{plim}}\widehat{\theta}=\theta
\]

Condição suficiente: tanto o viés $\lim_{n\rightarrow\infty}E\left(\widehat{\theta}_{n}\right)=\theta$ quanto a variância $\lim_{n\rightarrow\infty}var\left(\widehat{\theta}_{n}\right)=0$ devem tender a zero a medida que o tamanho da amostra cresce. Ou seja, o MSE deve tender a zero a medida que $n$ cresce.

Algumas propriedades dignas de nota:
\begin{enumerate}
\item Invariância (propriedade de Slutsky): Se $\widehat{\theta}$ for consistente e $h\left(\widehat{\theta}\right)$ for qualquer função que dependa apenas de $\widehat{\theta}$: $\underset{n\rightarrow\infty}{\text{plim}}h\left(\widehat{\theta}\right)=h\left(\theta\right)$
\item Se $b$ é uma constante: $\underset{n\rightarrow\infty}{\text{plim}}b=b$
\item Se $\widehat{\theta}_{1}$ e $\widehat{\theta}_{2}$ forem estimadores consistemtes:

\begin{align*}
\underset{n\rightarrow\infty}{\text{plim}}\left(\widehat{\theta}_{1}+\widehat{\theta}_{2}\right) & =\underset{n\rightarrow\infty}{\text{plim}}\left(\widehat{\theta}_{1}\right)+\underset{n\rightarrow\infty}{\text{plim}}\left(\widehat{\theta}_{2}\right)\\
\underset{n\rightarrow\infty}{\text{plim}}\left(\widehat{\theta}_{1}\widehat{\theta}_{2}\right) & =\underset{n\rightarrow\infty}{\text{plim}}\left(\widehat{\theta}_{1}\right)\underset{n\rightarrow\infty}{\text{plim}}\left(\widehat{\theta}_{2}\right)\\
\underset{n\rightarrow\infty}{\text{plim}}\left(\frac{\widehat{\theta}_{1}}{\widehat{\theta}_{2}}\right) & =\frac{\underset{n\rightarrow\infty}{\text{plim}}\left(\widehat{\theta}_{1}\right)}{\underset{n\rightarrow\infty}{\text{plim}}\left(\widehat{\theta}_{2}\right)}
\end{align*}

\item \textbf{Assintoticamente eficiente}: A variância da distribuição assintótica (quando $n\rightarrow\infty$) é chamada de variância assintótica de $\widehat{\theta}$. Então se $\widehat{\theta}$ for consistente e sua variância assintótica for menor do que a variância assintótica de todos os estimadores consistentes de $\theta$ então $\widehat{\theta}$ é chamado de assintoticamente eficiente.
\item Normalidade assintótica: o estimador $\widehat{\theta}$ é considerado ter distribuição assintoticamente normal se sua distribuição amostral tende a aproximar-se da distribuição normal à medida que o tamanho da amostra cresce.
\end{enumerate}
\textbf{Inferência estatística: estando as hipóteses}

O problema do teste de hipótese pode ser estabelecido da seguinte forma: admita que tenhamos uma variável aleatória $X$ com uma FDP conhecida f$\left(x,\theta\right)$, em que $\theta$ é o parâmetro da distribuição. Ao obtermos uma amostra aleatória de tamanho $n$, obtemos o estimador pontual $\widehat{\theta}$. 

Uma vez que o verdadeiro $\theta$ é raramente conhecido, levantamos a questão: o valor observado do estimador, obtido a partir da amostra, é plausível sob a distribuição que ele teria caso $\theta=\theta^{*}$. Isto é, a amostra poderia ter sido gerada por uma distribuição cujo parâmetro verdadeiro é $\theta^{*}$?
\begin{itemize}
\item \textbf{Hipótese nula (ou sustentada) $H_{0}$}: Na linguagem de teste da hipótese, ela afirma que os dados observados são compatíveis com a distribuição que o estimador $\widehat{\theta}$ teria caso o verdadeiro valor do parâmetro fosse $\theta=\theta^{*}$. 
\item \textbf{Hipótese alternativa} $H_{1}$ : A hipótese nula deve ser testada contra uma hipótese alternativa, por exemplo, $H_{1}$ $\theta\neq\theta^{*}$, ao qual afirma que o verdadeiro valor do parâmetro é diferente de $\theta^{*}$. Neste caso, a amostra observada não faz muito sentido se o parâmetro verdadeiro for $\theta^{*}$.
\end{itemize}
A hipótese nula e a alternativa podem ser simples ou compostas. Hipótese simples é quando se especifica os valor dos parâmetros da distribuição (ex.: $\sigma=2$), e hipóetse composta quando não (ex.: $\sigma>2$).

Para testar a hipótese nula, utilizamos a informação da amostra para obter o que é chamado de estatística de teste e então utilizamos a abordagem do intervalo de confiança ou do teste de significância para testar a hipótese nula. Camos trabalhar com um exemplo. Vamos supor que temos uma população de homens com altura média de $\mu$ e desvio padrão $\sigma=2.5$. Uma amostra de 100 homens foi tirada dessa população com média de $\overline{X}=67$. 

Ou seja:

\begin{align*}
X\sim N\left(\theta,\sigma^{2}\right) & =N\left(\mu,2.5^{2}\right)\\
\overline{X}=67 & n=100
\end{align*}

E vamos admitir que:

\begin{align*}
H_{0}: & \mu=\mu^{*}=69\\
H_{1}: & \mu\neq69
\end{align*}

A questão é: poderia a amostra com $\overline{X}=67$ (a estatística de teste), ter sido extraída da população com o valor médio de $\mu=69$?

\textbf{Abordagem do intervalo de confiança}

Tendo então a variável aleatória $X\sim N\left(\theta,\sigma^{2}\right)$ e tendo a estatística de teste $\overline{X}\sim\left(\mu,\sigma^{2}/n\right)$. Podemos então estabelecer um intervalo de confiança $\left(100-\alpha\right)$ para $\mu$, e verificar se esse intervalo inclui $\mu=\mu^{*}$. Se incluir não podemos rejeitar a hipótese nula, se não incluir, podemos rejeitar. Assim se $\alpha=0.05$ temos um intervalo de confiança de $95\%$ e se esse intervalo incluir $\mu^{*}$ não podemos rejeitar a hipótese nula e esperamos que 95 entre 100 intervalos assim estabelecidos deverão incluir $\mu^{*}$\footnote{Mas ainda há um risco de que no específico caso que realizarmos este teste, seja um dos 5\% que não incluiria $\mu^{*}.$}.

Usando os valores conhecidos de $\sigma$ e $n$$\overline{X}\sim N\left(\mu,\sigma^{2}/n\right)=N\left(\mu,\left(0.25\right)^{2}\right)$, Ou seja, $\sigma_{\overline{X}}=0.25$ é o erro padrão de $\overline{X}$. Como o intervalo entre $\overline{X}\pm1.96\sigma_{\overline{X}}$ cobre aproximadamente 95\% da area sob a normal, então o intervalo de confiança de 95\% de que $\mu$ esteja no intervalo é $\left[66.51,67.49\right]$. Uma vez que estabelecemos o intervalo, tudo que temos que fazer é verificar se $\mu$ está nesse intervalo. Se estiver não podemos rejeitara hipótese nula, se não estiver, podemos rejeitar. Como este intervalo não inclui $\mu^{*}=69$ então rejeitamos a hipótese nula com um coeficiente de confiança de 95\%\footnote{Mas observe que neste caso se a hipótese for que $\mu=\overline{X}$, nunca rejeitaríamos a hipótese. A hipótese então deve ser formulada independentemente dos dados observados.}.

Na linguagem do teste de hipóteses, o intervalo de confiança é chamada de região de aceitação, e as regiões fora de região crítica (ou região de rejeição). Os limites inferiores e superiores são chamados de valores críticos.

Temos ainda dois tipos de erros:
\begin{itemize}
\item \textbf{Erro do tipo 1}: rejeitar $H_{0}$ quando ela é verdadeira.
\item \textbf{Erro do tipo 2:} Não rejeitar $H_{0}$quando ela é falsa.
\end{itemize}
\textbf{Tipos de erro}

A abordagem clássica é supor que um erro do tipo 1 seja seja mais sério do que um erro do tipo 2. Ou seja primeiro se minimiza o erro do tipo 1, e então minimiza um erro do tipo 2 quando possível.

Temos:
\begin{itemize}
\item $\alpha$: \textbf{nível de significância} é a probabilidade de um erro do tipo 1.
\item $\beta$ é a probabilidade de um erro do tipo 2.
\begin{itemize}
\item Potência de teste: a probabilidade de não cometer um erro do tipo 2. Capacidade de rejeitar uma falsa hipótese nula.
\end{itemize}
\end{itemize}
Ou seja fixamos $\alpha$ e nívels como $0.01$ ou $0.05$ e maximizamos a potência de teste. Vamos tomar novamente um exemplo.

Vamos tomar $X\sim\left(\mu,100\right)$, uma amostra com 25 observações e considerar uma hipótese $H_{0}:\mu=50$. Temos também então $\overline{X}\sim N\left(\mu,100/25\right)$. Considerando um estimador com essa variância que tivesse média de $50$, a região de confiança seria $\pm1.6\sqrt{100/25}$ em torno da média, ou seja entre $46.08$ e $53.2$, rejeitamos então a hipótese nula de que a média é verdadeira se o valor da média da amostra estiver fora desse intervalo. Lembrando que obtemos isso a partir de que:

\[
P\left(\overline{X}-a\leq\mu\leq\overline{X}+a\right)=P\left(\mu-a\leq\overline{X}\leq\mu+a\right)
\]

Porém queremos saber a probabilidade de que $X$ esteja situado na região crítica se $\mu_{v}$ verdadeiro possuir um valor diferente de $50$ Vamos considerar três hitpóeses alternativas: $48,52,56.$ Nesse caso, considerando que $X$ tem 95\% de cair no intervalo $\left[46.08,53.2\right]$ de acordo com nossa hipótese, tem 5\% de chance de cair fora do intervalo e ser rejeitado se $\mu_{v}=50$. Podemos ver que caso a hipótese nula esteja correta, $\alpha$ nos a probabilidade de rejeitar a hipótese nula mesmo ela sendo verdadeira (erro do tipo I).

Aqui devemos nos atentar ao fato de que, em muitos testes estatísticos implementados em softwares como o Stata (por exemplo, o teste de correlação de Pearson), formulamos uma hipótese nula que representa ausência de efeito ou relação. No caso do teste de Pearson, a hipótese nula usual é a ausência de correlação linear na população, portanto quanto maior a probabilidade de que os dados observados serem compatíveis com a hipótese nula, mais fraco estatisticamente consideramos o resultado do nosso teste. Assim, quanto menor o nível de significância menos propensos estamos de cometer um erro do tipo I, isto é, descartar a hipótese nula mesmo ela sendo verdadeira. Se adotamos um valor muito baixo de $\alpha$ e mesmo assim rejeitamos a hipótese nula, supomos que é porque ela provavelmente é realmente falsa\footnote{Ainda que nunca possa ser $\alpha=0$, pois isso implicaria em nunca rejeitar a hipótese nula para nunca cometer nenhum erro do tipo I.}.

Mas se $\mu_{v}=48$ então sobre para 17\% a chance de rejeitar a hipótese nula , também temos 17\% para caso seja $\mu_{v}=52$ e 85\% caso seja $\mu_{v}=0.56$. Mas veja bem, se fosse $\mu=48$, então em cada 100 execuções, 83 vezes nossa média da amostra retornaria um um valor dentro do intervalo, ou seja, eu não rejeitaria a hipótese nula, mesmo estando errada e cometeríamos um erro do tipo II, assim, não rejeitar a hipótese nula não significa demonstrar que ela seja verdadeira; significa apenas que os dados observados não foram suficientemente improváveis sob a hipótese considerada. Da mesma forma, rejeitar a hipótese nula não significa demonstrar que ela seja falsa, mas sim que os dados observados foram suficientemente improváveis sob essa hipótese segundo o critério de significância adotado.

Podemos notar que pela natureza aleatória do experimento, conforme diminuimos $\alpha$, aumentamos a região de aceitação e podemos acabar aumentando a quantidade de erro do tipo II. Agora, vamos resumir os passos envolvidos no teste da hipótese estatística:
\begin{itemize}
\item \textbf{Passo 1}. Formule a hipótese nula $H_{0}$ e a hipótese alternativa $H_{1}$ (por exemplo: $H_{0}:\mu=69$ e $H_{1}:\mu\neq69$). 
\item \textbf{Passo 2}. Selecione a estatística de teste (por exemplo: $\overline{X}$ ).
\item \textbf{Passo 3}. Determine a distribuição de probabilidade da estatística de teste (por exemplo: $\overline{X}\sim N\left(\mu,\sigma^{2}/n\right)$. 
\item \textbf{Passo 4}. Escolha o nível de significância $\alpha$ (a probabilidade de cometer um erro do tipo I: rejeitar $H_{0}$ quando ela é verdadeira).
\item \textbf{Passo 5}. Utilizando a distribuição de probabilidade da estatística de teste, estabeleça um valor de confiança$100(1-\alpha)\%$. Se o valor hipotetizado em $H_{0}$ (por exemplo: $\mu=69$) pertencer ao intervalo de confiança, não rejeitamos $H_{0}$ ao nível $\alpha$. Porém, se ele estiver fora desse intervalo (ou seja, dentro da região de rejeição), pode-se rejeitar a hipótese nula. Tenha em mente que, ao rejeitar uma hipótese nula, corre-se o risco de estar errado em uma porcentagem $\alpha$, e ao não rejeitar também temos uma chance $\beta$ de estarmos errados.
\end{itemize}

\textbf{Valor p ou nível exato de significância}

Existe uma alternativa mais informativa do que simplesmente pré-selecionar níveis arbitrários de significância $\alpha$. Além disso, os exemplos anteriores assumiam que a variância da população era conhecida, hipótese frequentemente irrealista, já que na prática normalmente não conhecemos nem a média nem a variância populacional.

O valor-p (ou nível exato de significância) de uma estatística de teste é definido como o menor nível de significância ao qual a hipótese nula pode ser rejeitada. Para entender essa ideia, suponha a hipótese nula $H_{0}:\mu=a.$ A partir de uma amostra, calculamos a média amostral $\overline{X}$e para um dado nível de significância $\alpha$, podemos construir um intervalo de confiança de nível$\left(1-\alpha\right)\%$. Se o valor hipotético $a$ estiver fora desse intervalo, rejeitamos $H_{0}$; caso contrário, não a rejeitamos.

Caso tivéssemos rejeitado a hipótese nula, poderíamos imaginar diminuir $\alpha$. Como isso aumenta o intervalo de confiança, torna-se mais difícil rejeitar $H_{0}$. Nesse contexto, o valor-$p$ representa exatamente o ponto limite: ele é o menor valor de $\alpha$ para o qual a hipótese nula ainda seria rejeitada pelos dados observados. Ou seja, quanto menor o valor-p, maior precisaria ser o intervalo de confiança para que $H_{0}$ deixasse de ser rejeitada. Equivalentemente, isso significa que os dados observados são menos compatíveis com a hipótese nula, pois seriam mais improváveis caso $H_{0}$ fosse verdadeira. Ou seja, para um valor-p de $5\%$ por exemplo, aassumindo que a hipótese nula seja verdadeira, existe menos de 5\% de probabilidade de observar um resultado tão extremo quanto o obtido.

Como mencionado, em grande parte dos testes estatísticos, a hipótese nula $H_{0}$ representa uma hipótese de ausência de efeito, ausência de diferença ou ausência de relação entre variáveis. Nesse contexto, o objetivo do teste é procurar evidências contra a hipótese nula. Assim, quanto menor o valor-p, menos compatíveis os dados observados parecem ser com a hipótese nula e, portanto, maior a evidência contra $H_{0}$. Entretanto, é importante ressaltar que o teste frequentista não prova a hipótese alternativa; ele apenas avalia quão incompatíveis os dados observados são com a hipótese nula.

\textbf{Abordagem do teste da significância}

Vamos abordar uma variável aleatória normal normalizada:

\[
Z=\frac{\overline{X}-\mu}{\sigma/\sqrt{n}}\sim N\left(0,1\right)
\]

$\overline{X}$ e $n$ são conhecidos, e se $\sigma$ for especificado e considerarmos um $\mu=\mu^{*}$específico então podemos calcular a probabilidade de ter um valor igual ou maior que $Z$. Se a probabilidade for muito pequena, então podemos rejeitar a hipótese nula, isto é, se a hipótese nula fosse verdadeira, deveríamos ter uma probabilidade alta de obter $Z$ ou maior. Voltando ao nosso exemplo, suponha a hipótese nula $H_{0}:\mu=69$. Nesse caso,

\[
Z=\frac{67-69}{2.5/\sqrt{100}}=-8
\]
.

A probabilidade de obter um valor de $Z$ tão extremo quanto $\left|8\right|$ (ou mais), assumindo que $H_{0}$ seja verdadeira, é extremamente pequena. Isso significa que, se de fato $\mu=69$, seria extremamente improvável observar uma média amostral tão distante quanto $\overline{X}=67$. Portanto, suspeitamos que nossa amostra não provenha de uma população cuja média verdadeira seja $69$.

Na linguagem dos testes de hipótese, dizemos que uma estatística de teste é estatisticamente significativa quando a probabilidade de observar um resultado tão extremo quanto o obtido é menor ou igual ao nível de significância $\alpha$, isto é, a probabilidade máxima tolerada de cometer um erro do tipo I (rejeitar $H_{0}$ quando ela é verdadeira). Análogo a discussão anterior, assumindo que a hipótese é verdadeira se escolhemos $\alpha=0.05$. temos um intervalo de confiança de 95\%, mas ainda sim, 5\% dos dados coletados aleatoriamente poderiam ter um valores extremos fora do intervalo e rejeitaríamos a hipótese nula mesmo sendo verdadeira. Se definirmos $\alpha=0.05$,então rejeitamos a hipótese nula sempre que a probabilidade associada ao valor observado da estatística (ou maior) de teste for menor que $5\%$. Em nosso exemplo, a probabilidade de obter um valor tão extremo quanto $Z=-8$ é muito menor que $2.5\%$ em cada cauda da distribuição normal. Portanto, consideramos o resultado estatisticamente significativo e rejeitamos a hipótese nula de que o verdadeiro valor de $\mu$ seja $69$. Na prática então esta abordagem nos dá p-value.

A estrutura geral é então:
\begin{itemize}
\item \textbf{Passo 1.} Formule a hipótese nula $H_{0}$ e a hipótese alternativa $H_{1}$
\item \textbf{Passo 2}. Selecione a estatística de teste.
\item \textbf{Passo 3}. Determine a distribuição de probabilidade da estatística de teste.
\item \textbf{Passo 4}.Calcular a probabilidade, assumindo $H_{0}$ verdadeira, de obter um resultado tão extremo quanto o observado.
\end{itemize}

\bibliographystyle{plain}
\bibliography{tese,Artigo1/references,Artigo2/referencias}

\end{document}
